\documentclass[10pt]{article}

\usepackage[T1]{fontenc}
\usepackage{amsmath,amssymb,amsthm}
\usepackage{array,booktabs,longtable}
\usepackage[colorlinks=true,linkcolor=blue,urlcolor=blue]{hyperref}
\usepackage[a4paper,margin=2cm]{geometry}

\setlength{\parindent}{14pt}
\setlength{\parskip}{2pt}
\setlength{\LTpre}{3pt}
\setlength{\LTpost}{3pt}
\allowdisplaybreaks

\newtheorem{proposition}{Proposition}
\renewcommand{\theproposition}{S\arabic{proposition}}
\renewcommand{\thesection}{S\arabic{section}}
\renewcommand{\thetable}{S\arabic{table}}

\begin{document}

\begin{center}
{\Large\bfseries Supplementary Proof Details}\\[4pt]
{\large for \emph{The complete equidistribution classification of\\
length-three mesh patterns on layered permutations}}\\[7pt]
Jiahao Zhang
\end{center}

\noindent
This note records the two finite ingredients used at the end of the
classification: the incidence data for the coordinate identifications and
the exact separation of the resulting components.  All identities valid for
arbitrary host size, including the signature formula, the pointwise
criterion, the kernel-degree theorem, the coordinate lemma, and the eighteen
distributional bridges, are proved in the main paper.

\section{The last two finite quotients}

After pointwise equality, reverse-complement, and the complete
kernel-degree reduction, the main paper obtains $200$ classes.  Of these,
$28$ are the unique-support classes.  Among the remaining $172$ classes,
$84$ are isolated under coordinate relabelling, while the other $88$ are
joined by the $43$ edges of Table~\ref{tab:coordinate-witnesses}.  Seven
edges close triangles already connected by two earlier edges, so the rank
drop is $36$ and
\[
28+84+(88-36)=164.
\]

The five all-size bridge families proved in the main paper give the inventory
in Table~\ref{tab:bridge-inventory}.

\begin{table}[htbp]
\centering
\small
\caption{The eighteen effective bridges from $164$ to $146$ components.}
\label{tab:bridge-inventory}
\begin{tabular}{@{}l c r@{}}
\toprule
mechanism & parameter range & effective joins\\
\midrule
uniform cut shift & $(a,b)\in\{0,1,2\}^2$ & $9$\\
exceptional rational identity & Class $79$ & $1$\\
non-overlapping contraction & $z\in\{0,1,2,3\}$ & $4$\\
separator complementation & $a\in\{0,1,2\}$ & $3$\\
MacMahon bridge & major index--inversion & $1$\\
\midrule
total && $18$\\
\bottomrule
\end{tabular}
\end{table}

For every listed bridge, its two endpoints lie in distinct coordinate
components, and no bridge becomes redundant after the preceding bridges are
added.  Thus all eighteen joins are effective and the resulting graph has
$164-18=146$ components.  This is the upper-bound partition used in the
complete-classification theorem.

\section{The 43 coordinate edges}

We use the notation $(r;z;E)$ from the main paper, omitting commas inside the
block word, internal-channel word, and external-set word.  Thus
$(111;012;03)$ denotes $r=(1,1,1)$, $z=(0,1,2)$, and $E=\{0,3\}$.

Every row of Table~\ref{tab:coordinate-witnesses} is oriented.  The appendix
of the main paper proves, symbolically for arbitrary layer number $m$ and
arbitrary positive layer sizes, the identity
\[
P_{\Sigma}(L)=P_{\Sigma'}(U_i(L))
\]
displayed by that row.  Thus the table records finite incidence data for
all-size mathematical identities; it is not a list of bounded numerical
tests.

{\footnotesize
\setlength{\tabcolsep}{4pt}
\begin{longtable}{@{}r c c p{10.3cm}@{}}
\caption{Explicit oriented coordinate edges.  An entry
$\Sigma\longrightarrow\Sigma'$ in the $U_i$ column means
$P_{\Sigma}(L)=P_{\Sigma'}(U_i(L))$ for every positive host composition
$L$.\label{tab:coordinate-witnesses}}\\
\toprule
class & local edge & map & oriented signature identity\\
\midrule
\endfirsthead
\multicolumn{4}{c}{\tablename\ \thetable\ continued}\\
\toprule
class & local edge & map & oriented signature identity\\
\midrule
\endhead
\midrule
\multicolumn{4}{r}{continued on next page}\\
\endfoot
\bottomrule
\endlastfoot
3 & $1$--$2$ & $U_1$ & $(111;000;01)\longrightarrow(111;000;03)$\\
11 & $1$--$2$ & $U_1$ & $(111;001;01)\longrightarrow(111;010;03)$\\
20 & $1$--$2$ & $U_1$ & $(111;002;01)\longrightarrow(111;020;03)$\\
12 & $1$--$2$ & $U_2$ & $(111;001;02)\longrightarrow(111;001;23)$\\
12 & $1$--$3$ & $U_3$ & $(111;001;02)\longrightarrow(111;010;01)$\\
12 & $2$--$3$ & $U_4$ & $(111;001;23)\longrightarrow(111;010;01)$\\
22 & $1$--$2$ & $U_2$ & $(111;002;03)\longrightarrow(111;002;23)$\\
22 & $1$--$3$ & $U_3$ & $(111;002;03)\longrightarrow(111;020;01)$\\
22 & $2$--$3$ & $U_4$ & $(111;002;23)\longrightarrow(111;020;01)$\\
43 & $1$--$2$ & $U_2$ & $(111;012;03)\longrightarrow(111;102;23)$\\
43 & $1$--$3$ & $U_3$ & $(111;012;03)\longrightarrow(111;021;01)$\\
43 & $2$--$3$ & $U_4$ & $(111;102;23)\longrightarrow(111;021;01)$\\
13 & $1$--$2$ & $U_8$ & $(111;001;0)\longrightarrow(111;010;0)$\\
21 & $1$--$2$ & $U_8$ & $(111;002;02)\longrightarrow(111;020;02)$\\
23 & $1$--$2$ & $U_8$ & $(111;002;0)\longrightarrow(111;020;0)$\\
44 & $1$--$2$ & $U_8$ & $(111;012;0)\longrightarrow(111;021;0)$\\
35 & $1$--$2$ & $U_5$ & $(111;011;01)\longrightarrow(111;101;01)$\\
37 & $1$--$2$ & $U_6$ & $(111;011;12)\longrightarrow(111;101;02)$\\
38 & $1$--$2$ & $U_6$ & $(111;011;1)\longrightarrow(111;101;0)$\\
65 & $1$--$2$ & $U_6$ & $(111;022;23)\longrightarrow(111;202;03)$\\
90 & $1$--$2$ & $U_6$ & $(111;122;23)\longrightarrow(111;212;03)$\\
102 & $1$--$2$ & $U_6$ & $(111;222;23)\longrightarrow(111;222;03)$\\
42 & $1$--$2$ & $U_5$ & $(111;012;01)\longrightarrow(111;102;01)$\\
42 & $1$--$3$ & $U_1$ & $(111;012;01)\longrightarrow(111;021;02)$\\
42 & $2$--$3$ & $U_7$ & $(111;102;01)\longrightarrow(111;021;02)$\\
46 & $1$--$2$ & $U_6$ & $(111;012;13)\longrightarrow(111;102;03)$\\
46 & $1$--$3$ & $U_{11}$ & $(111;012;13)\longrightarrow(111;021;23)$\\
46 & $2$--$3$ & $U_7$ & $(111;102;03)\longrightarrow(111;021;23)$\\
48 & $1$--$2$ & $U_9$ & $(111;012;3)\longrightarrow(111;102;3)$\\
63 & $1$--$2$ & $U_9$ & $(111;022;13)\longrightarrow(111;202;13)$\\
67 & $1$--$2$ & $U_9$ & $(111;022;3)\longrightarrow(111;202;3)$\\
92 & $1$--$2$ & $U_9$ & $(111;122;3)\longrightarrow(111;212;3)$\\
53 & $1$--$2$ & $U_7$ & $(111;102;02)\longrightarrow(111;021;12)$\\
55 & $1$--$2$ & $U_7$ & $(111;102;0)\longrightarrow(111;021;2)$\\
59 & $1$--$2$ & $U_3$ & $(111;022;03)\longrightarrow(111;022;01)$\\
59 & $1$--$3$ & $U_2$ & $(111;022;03)\longrightarrow(111;202;23)$\\
59 & $2$--$3$ & $U_{10}$ & $(111;022;01)\longrightarrow(111;202;23)$\\
87 & $1$--$2$ & $U_3$ & $(111;122;03)\longrightarrow(111;122;01)$\\
87 & $1$--$3$ & $U_2$ & $(111;122;03)\longrightarrow(111;212;13)$\\
87 & $2$--$3$ & $U_{10}$ & $(111;122;01)\longrightarrow(111;212;13)$\\
79 & $1$--$2$ & $U_1$ & $(111;112;01)\longrightarrow(111;121;02)$\\
81 & $1$--$2$ & $U_1$ & $(111;112;0)\longrightarrow(111;121;2)$\\
80 & $1$--$2$ & $U_{11}$ & $(111;112;03)\longrightarrow(111;121;23)$\\
\end{longtable}

}

\section{Finite separation of the 146 components}

\begin{proposition}[Pairwise separation]\label{prop:finite-separation}
Choose one signature from each of the $146$ components obtained in
Sections~S1--S2.  Representatives can be chosen so that every two distinct
representatives have different distribution polynomials at some host size
$n\le7$.  Among the $\binom{146}{2}=10{,}585$ unordered pairs, the numbers
whose first separating sizes are $4,5,6,7$ are, respectively,
\[
10{,}260,\qquad302,\qquad21,\qquad2.
\]
Consequently the $146$ components are pairwise distribution-inequivalent.
\end{proposition}

\begin{proof}
\enlargethispage{2\baselineskip}
For a signature $\Sigma$, substitute the occurrence formula from the main
paper into
\[
D_{\Sigma,n}(q)
=\sum_{m=1}^{n}\ \sum_{\substack{\ell_1+\cdots+\ell_m=n\\
\ell_i\ge1}}q^{P_\Sigma(\ell_1,\ldots,\ell_m)}.
\]
For $n\le7$ this is a finite sum of monomials with integer coefficients.
Exact coefficient comparison gives the four numbers in the statement.
The two pairs that require $n=7$ are displayed below; all other pairs
separate by $n\le6$.

Put
\[
\Sigma_A=(111;000;\{0,2\}),\quad
\Sigma_B=(111;000;\{1,2\}).
\]
Then
\[
\begin{aligned}
D_{\Sigma_A,7}(q)
 &=7+13q^5+5q^6+4q^7+18q^8+8q^9+2q^{10}+7q^{12},\\
D_{\Sigma_B,7}(q)
 &=7+13q^5+4q^6+6q^7+17q^8+8q^9+2q^{10}+7q^{12}.
\end{aligned}
\]
For
\[
\Sigma_C=(111;222;\{0,2\}),\quad
\Sigma_D=(111;222;\{1,2\}),
\]
one has
\[
\begin{aligned}
D_{\Sigma_C,7}(q)&=49+7q+5q^2+2q^3+q^5,\\
D_{\Sigma_D,7}(q)&=48+9q+4q^2+2q^3+q^5.
\end{aligned}
\]
A difference at a single size rules out equidistribution.  Therefore no two
of the $146$ components can merge.
\end{proof}

\end{document}
